\chapter{Classification and Kinship}\label{ch:g6:classification-kinship}

Why does the \omterm{def:g4:classifying-animals:classification}{classification} of
\cref{ch:g4:classifying-animals} work so well? Sort \omterm{def:g1:animals-around-us:animal}{animals}
by \omterm{def:g1:animals-around-us:covering}{feathers} and \omterm{def:g1:animals-around-us:covering}{fur}, and a thousand other features fall
obediently into line --- milk-feeders have \omterm{def:g4:classifying-animals:vertebrate}{backbones},
feather-wearers lay shelled \omterm{def:g2:animals-grow-up:born}{eggs}, and no one has to check.
\Cref{rem:g4:classifying-animals:tree} left the question
hanging. This chapter answers it with one word --- \omterm{prop:g6:classification-kinship:kinship}{kinship}
--- and rebuilds the boxes as a family \omterm{def:g6:classification-kinship:tree}{tree}.

\section{Attributes, shared and nested}

\begin{definition}[Attribute]\label{def:g6:classification-kinship:attribute}
An \emph{attribute}\index{attribute} is a feature an
organism \emph{possesses}, stated so it can be checked: has
a \omterm{def:g3:bones-and-muscles:skeleton}{skeleton} with a \omterm{def:g4:classifying-animals:vertebrate}{backbone}; has four \omterm{def:g1:my-body:parts}{limbs}; has \omterm{def:g1:animals-around-us:covering}{fur} and
milk-feeding; has \omterm{def:g1:animals-around-us:covering}{feathers}. \omterm{def:g4:classifying-animals:classification}{Classification} is built on
attributes --- \cref{def:g3:sorting-living-things:criterion}'s
good criteria, formalized.
\end{definition}

\begin{proposition}[Attributes nest]\label{prop:g6:classification-kinship:nest}
Checked across many \omterm{def:g5:biodiversity:species}{species}, \omterm{def:g6:classification-kinship:attribute}{attributes} fall into
\emph{nested} sets --- every \omterm{def:g1:animals-around-us:animal}{animal} with \omterm{def:g1:animals-around-us:covering}{fur} also has four
\omterm{def:g1:my-body:parts}{limbs}; every four-limbed \omterm{def:g1:animals-around-us:animal}{animal} also has a \omterm{def:g4:classifying-animals:vertebrate}{backbone}; never
the reverse pattern (\omterm{def:g1:animals-around-us:covering}{fur} without \omterm{def:g4:classifying-animals:vertebrate}{backbone}) and never
crossing sets (\omterm{def:g1:animals-around-us:covering}{fur} and \omterm{def:g1:animals-around-us:covering}{feathers} together). That nesting is
why groups-inside-groups
(\cref{def:g4:classifying-animals:classification}) fits the
living world at all.
\end{proposition}
\admitted

\begin{omfigure}
\begin{tikzpicture}[scale=0.9]
  % nested attribute boxes
  \draw[thick, omDef, fill=omDef!4] (0,0) rectangle (11.4,4.4);
  \node[font=\small\bfseries, omDef] at (2.2,4.1)
    {has a backbone};
  \draw[thick, omMeth, fill=omMeth!5] (0.4,0.35) rectangle (8.6,3.5);
  \node[font=\small\bfseries, omMeth] at (2.35,3.2)
    {has four limbs};
  \draw[thick, omProp, fill=omProp!6] (0.8,0.7) rectangle (5.6,2.6);
  \node[font=\small\bfseries, omProp] at (2.6,2.3)
    {has fur and milk};
  \node[font=\small] at (3.2,1.4) {cat, bat, whale};
  \node[font=\small, align=center] at (7.0,1.6)
    {frog, lizard,\\ blackbird};
  \node[font=\small] at (10.0,1.6) {trout};
\end{tikzpicture}

{\small \omterm{def:g6:classification-kinship:attribute}{Attributes} nest: fur-bearers sit inside the
four-limbed, the four-limbed inside the backboned. The trout
has the \omterm{def:g4:classifying-animals:vertebrate}{backbone} but stands outside the four-limbed box.}
\end{omfigure}

\begin{example}[Placing three animals]\label{ex:g6:classification-kinship:placing}
Check the boxes for the whale, the trout and the blackbird.
Whale: \omterm{def:g4:classifying-animals:vertebrate}{backbone} yes, four \omterm{def:g1:my-body:parts}{limbs} yes (front flippers, and
hidden vestiges behind), fur-and-milk yes --- innermost box,
with the cat. Trout: \omterm{def:g4:classifying-animals:vertebrate}{backbone} yes, stop --- outer box only.
Blackbird: \omterm{def:g4:classifying-animals:vertebrate}{backbone}, four \omterm{def:g1:my-body:parts}{limbs} (two \omterm{def:g1:my-body:parts}{legs}, two wings),
\omterm{def:g1:animals-around-us:covering}{feathers} --- the four-limbed box, in \omterm{def:g1:animals-around-us:covering}{feathers}' own
compartment. No \omterm{def:g1:animals-around-us:animal}{animal} ever demands two crossing boxes at
once.
\end{example}

\section{Kinship: the reason behind the nesting}

\begin{proposition}[Shared attributes mean shared ancestry]\label{prop:g6:classification-kinship:kinship}
The explanation of the nesting is inheritance
(\cref{rem:g4:animal-reproduction:mix}: young receive from
their parents). \omterm{def:g5:biodiversity:species}{Species} sharing an \omterm{def:g6:classification-kinship:attribute}{attribute} share it
because they \emph{inherited it from a common ancestor} ---
a distant grandparent \omterm{def:g5:biodiversity:species}{species}. The more \omterm{def:g6:classification-kinship:attribute}{attributes} two
\omterm{def:g5:biodiversity:species}{species} share, the closer their
\emph{kinship}\index{kinship}: the cat and the whale, both
furred milk-feeders, are closer cousins to each other than
either is to the trout. A \omterm{def:g4:classifying-animals:classification}{classification} group is, at
bottom, a \emph{family}: an ancestor and all its
descendants.
\end{proposition}
\admitted

\begin{example}[Reading the whale correctly at last]\label{ex:g6:classification-kinship:whale}
This is why the whale puzzle
(\cref{exo:g4:classifying-animals:9}) had to resolve as it
did: fur-and-milk marks the whale's family line --- its
ancestors were milk-feeding land \omterm{prop:g3:sorting-living-things:groups}{mammals} --- while its
fish-like shape is the water's doing
(\cref{exo:g4:adapted-to-their-world:11}), stamped on an
unrelated family by the same way of life. \omterm{def:g6:classification-kinship:attribute}{Attributes} of
ancestry sort families; resemblances of \omterm{def:g2:where-animals-live:habitat}{habitat} can be
borrowed by strangers. \omterm{def:g4:classifying-animals:classification}{Classification} trusts the first and
distrusts the second.
\end{example}

\begin{definition}[The tree of kinship]\label{def:g6:classification-kinship:tree}
\omterm{prop:g6:classification-kinship:kinship}{Kinship} is drawn as a \emph{tree}\index{tree of kinship}:
\omterm{def:g1:plants-around-us:tree}{branch} points stand for common ancestors, and each
\omterm{def:g6:classification-kinship:attribute}{attribute} appears once, at the \omterm{def:g1:plants-around-us:tree}{branch} where it arose ---
every \omterm{def:g5:biodiversity:species}{species} beyond that fork inherits it. The nested
boxes and the tree carry the same information: a box is a
\omterm{def:g1:plants-around-us:tree}{branch} with all its twigs.
\end{definition}

\begin{omfigure}
\begin{tikzpicture}[scale=0.95]
  % tree
  \draw[very thick] (0,0) -- (2.2,0);
  % backbone mark
  \fill[omDef] (1.6,0) circle (0.09);
  \node[font=\footnotesize, omDef, align=center] at (1.6,-0.55)
    {backbone\\ appears};
  \draw[very thick] (2.2,0) -- (3.6,1.6) -- (10.4,1.6);
  \draw[very thick] (2.2,0) -- (4.4,-1.4) -- (10.4,-1.4);
  \node[font=\small, right] at (10.5,-1.4) {trout};
  % four limbs mark
  \fill[omMeth] (3.1,1.05) circle (0.09);
  \node[font=\footnotesize, omMeth, align=center] at (2.4,1.45)
    {four limbs\\ appear};
  \draw[very thick] (5.0,1.6) -- (6.2,2.8) -- (10.4,2.8);
  % fur mark on upper branch
  \fill[omProp] (5.75,2.35) circle (0.09);
  \node[font=\footnotesize, omProp, align=center] at (5.0,2.9)
    {fur and milk\\ appear};
  \draw[very thick] (6.9,2.8) -- (7.7,3.7) -- (10.4,3.7);
  \node[font=\small, right] at (10.5,3.7) {cat};
  \node[font=\small, right] at (10.5,2.8) {whale};
  % feathers on middle branch
  \fill[omThm] (7.4,1.6) circle (0.09);
  \node[font=\footnotesize, omThm, align=center] at (7.4,1.05)
    {feathers\\ appear};
  \node[font=\small, right] at (10.5,1.6) {blackbird};
\end{tikzpicture}

{\small A pocket \omterm{def:g6:classification-kinship:tree}{tree of kinship}. Each colored dot is an
\omterm{def:g6:classification-kinship:attribute}{attribute} appearing once; every \omterm{def:g1:plants-around-us:tree}{branch} beyond it inherits.
Cat and whale share the newest fork --- the closest cousins
drawn here.}
\end{omfigure}

\begin{method}[Building a small kinship tree]\label{met:g6:classification-kinship:build}
Given a handful of \omterm{def:g5:biodiversity:species}{species}:
\begin{enumerate}
  \item table their \omterm{def:g6:classification-kinship:attribute}{attributes} --- \omterm{def:g5:biodiversity:species}{species} in rows, checkable
        \omterm{def:g6:classification-kinship:attribute}{attributes} in columns, ticks where possessed;
  \item group by shared ticks, widest \omterm{def:g6:classification-kinship:attribute}{attribute} first: the
        nesting appears in the table;
  \item draw the \omterm{def:g6:classification-kinship:tree}{tree}: one fork per \omterm{def:g6:classification-kinship:attribute}{attribute}, the widest
        nearest the trunk; place each \omterm{def:g5:biodiversity:species}{species} past exactly
        the forks whose \omterm{def:g6:classification-kinship:attribute}{attributes} it possesses;
  \item read \omterm{prop:g6:classification-kinship:kinship}{kinship}: two \omterm{def:g5:biodiversity:species}{species}' closeness is the last
        fork they share.
\end{enumerate}
\end{method}

\begin{remark}[What the tree quietly implies]\label{rem:g6:classification-kinship:implies}
Take the \omterm{def:g6:classification-kinship:tree}{tree} seriously and it says something enormous:
follow any two \omterm{def:g1:plants-around-us:tree}{branches} back and they \emph{meet}. Cat and
trout share a grandparent \omterm{def:g5:biodiversity:species}{species}; so, further back, must
cat and snail, daisy and child. All \omterm{def:g6:exploring-our-environment:components}{living things} would then
be one family, its unity
(\cref{prop:g5:first-look-at-cells:principle}) a family
resemblance --- and \omterm{def:g5:biodiversity:species}{species} must be able to \emph{change}
across generations, or the branching could never have
happened. Whether that is true, and how it could work, is
the great question of \cref{ch:g9:evidence-of-evolution}'s
year --- this \omterm{def:g6:classification-kinship:tree}{tree} is the diagram that forces it.
\end{remark}

\section{Exercises}

\begin{exercise}[$\star$]\label{exo:g6:classification-kinship:1}
What is an \omterm{def:g6:classification-kinship:attribute}{attribute}? Give three, stated checkably.
\end{exercise}

\begin{exercise}[$\star$]\label{exo:g6:classification-kinship:2}
What does it mean that \omterm{def:g6:classification-kinship:attribute}{attributes} \emph{nest}? Give the
chapter's three-box example.
\end{exercise}

\begin{exercise}[$\star$]\label{exo:g6:classification-kinship:3}
Why do \omterm{def:g5:biodiversity:species}{species} share \omterm{def:g6:classification-kinship:attribute}{attributes}, by
\cref{prop:g6:classification-kinship:kinship}?
\end{exercise}

\begin{exercise}[$\star$]\label{exo:g6:classification-kinship:4}
On the pocket \omterm{def:g6:classification-kinship:tree}{tree}, which \omterm{def:g6:classification-kinship:attribute}{attribute} dot does the trout
stand beyond? Which dots does the whale stand beyond?
\end{exercise}

\begin{exercise}[$\star$]\label{exo:g6:classification-kinship:5}
Who are closer cousins: whale and cat, or whale and trout?
Read the answer from the \omterm{def:g6:classification-kinship:tree}{tree}'s forks.
\end{exercise}

\begin{exercise}[$\star$]\label{exo:g6:classification-kinship:6}
State the relation between a nested box and a \omterm{def:g6:classification-kinship:tree}{tree} \omterm{def:g1:plants-around-us:tree}{branch}.
\end{exercise}

\begin{exercise}[$\star\star$]\label{exo:g6:classification-kinship:7}
Why does \omterm{def:g4:classifying-animals:classification}{classification} trust fur-and-milk but distrust
``shaped like a \omterm{prop:g3:sorting-living-things:groups}{fish}''? Use
\cref{ex:g6:classification-kinship:whale}.
\end{exercise}

\begin{exercise}[$\star\star$]\label{exo:g6:classification-kinship:8}
Build the \omterm{def:g6:classification-kinship:attribute}{attribute} table of
\cref{met:g6:classification-kinship:build} for: trout,
frog, cat, blackbird --- \omterm{def:g6:classification-kinship:attribute}{attributes}: \omterm{def:g4:classifying-animals:vertebrate}{backbone}, four \omterm{def:g1:my-body:parts}{limbs},
fur-and-milk, \omterm{def:g1:animals-around-us:covering}{feathers}.
\end{exercise}

\begin{exercise}[$\star\star$]\label{exo:g6:classification-kinship:9}
Add the lizard to the pocket \omterm{def:g6:classification-kinship:tree}{tree}: which forks does it
pass, and where does its \omterm{def:g1:plants-around-us:tree}{branch} leave?
\end{exercise}

\begin{exercise}[$\star\star$]\label{exo:g6:classification-kinship:10}
The bat flies like a blackbird but nurses like a cat.
Which resemblance marks \omterm{prop:g6:classification-kinship:kinship}{kinship} and which is borrowed by
way of life? Decide with the \omterm{def:g6:classification-kinship:tree}{tree}'s logic.
\end{exercise}

\begin{exercise}[$\star\star$]\label{exo:g6:classification-kinship:11}
``A group is an ancestor and all its descendants.'' Test
the sentence on the \omterm{prop:g3:sorting-living-things:groups}{mammals} and on
\cref{exo:g4:classifying-animals:11}'s platypus: does the
egg-laying oddity break the family, or belong to it?
\end{exercise}

\begin{exercise}[$\star\star\star$]\label{exo:g6:classification-kinship:12}
Follow \cref{rem:g6:classification-kinship:implies} to its
edge: what must be true of \omterm{def:g5:biodiversity:species}{species} over long time, if
every pair of \omterm{def:g1:plants-around-us:tree}{branches} really meets? State the implication
in one sentence --- and name what kind of evidence you
would demand for it.
\end{exercise}

\section{Problem: The Museum Drawer}

\begin{problem}[{Weekend problem --- six specimens, one
attribute table, one tree}]\label{pb:g6:classification-kinship:1}
A museum lends the class a drawer of six specimens: a
trout, a frog, a lizard, a pigeon, a mouse, and a
dolphin's \omterm{def:g3:bones-and-muscles:skeleton}{skull} with flipper casts. The task: classify
them by \omterm{def:g6:classification-kinship:attribute}{attributes} and draw their \omterm{prop:g6:classification-kinship:kinship}{kinship} \omterm{def:g6:classification-kinship:tree}{tree}.

\medskip\noindent\textbf{Part I --- The table.}
\begin{enumerate}
  \item Build the table for the six, with \omterm{def:g6:classification-kinship:attribute}{attributes}:
        \omterm{def:g4:classifying-animals:vertebrate}{backbone}; four \omterm{def:g1:my-body:parts}{limbs}; dry scaly \omterm{def:g1:the-five-senses:sense}{skin}; \omterm{def:g1:animals-around-us:covering}{feathers};
        \omterm{def:g1:animals-around-us:covering}{fur} and \omterm{def:g2:animals-grow-up:born}{milk}. (For the dolphin, near-hairless,
        the museum label vouches: milk-fed young ---
        \cref{ex:g3:sorting-living-things:surprises}
        explains why that testimony settles it.)
  \item Which \omterm{def:g6:classification-kinship:attribute}{attribute} do all six share? What does the
        widest sharing mean, in \omterm{prop:g6:classification-kinship:kinship}{kinship} terms?
  \item Which specimens tick ``four \omterm{def:g1:my-body:parts}{limbs}''? Explain the
        pigeon's and the dolphin's ticks --- what are
        wings and flippers, on the \omterm{def:g6:classification-kinship:tree}{tree}'s logic?
  \item Why does ``lives in water'' appear nowhere in the
        table, though trout and dolphin both do? Cite the
        chapter's rule.
\end{enumerate}

\medskip\noindent\textbf{Part II --- The \omterm{def:g6:classification-kinship:tree}{tree}.}
\begin{enumerate}[resume]
  \item Draw the \omterm{def:g6:classification-kinship:tree}{tree}: order the forks --- \omterm{def:g4:classifying-animals:vertebrate}{backbone}, four
        \omterm{def:g1:my-body:parts}{limbs}, then the coats (scales, \omterm{def:g1:animals-around-us:covering}{feathers}, \omterm{def:g1:animals-around-us:covering}{fur}) on
        their own \omterm{def:g1:plants-around-us:tree}{branches} --- and place all six.
  \item Read from your \omterm{def:g6:classification-kinship:tree}{tree}: the mouse's closest cousin
        in the drawer, and the trout's.
  \item The frog stands past ``four \omterm{def:g1:my-body:parts}{limbs}'' but before
        every coat fork. What does its bare, moist \omterm{def:g1:the-five-senses:sense}{skin}
        (\cref{rem:g4:classifying-animals:confusion})
        mark it as, and where does its \omterm{def:g1:plants-around-us:tree}{branch} leave?
  \item A label falls off a seventh specimen: dry scaly
        \omterm{def:g1:the-five-senses:sense}{skin}, four \omterm{def:g1:my-body:parts}{limbs}, \omterm{def:g4:classifying-animals:vertebrate}{backbone}. Place it on the \omterm{def:g6:classification-kinship:tree}{tree}
        and name its class --- without seeing the label.
\end{enumerate}

\medskip\noindent\textbf{Part III --- The meaning.}
\begin{enumerate}[resume]
  \item The dolphin and the trout share streamlining;
        the dolphin and the mouse share fur-and-milk
        ancestry. Which pair sits closer on the \omterm{def:g6:classification-kinship:tree}{tree}, and
        why must the \omterm{def:g6:classification-kinship:tree}{tree} ignore the streamlining?
  \item State what the fork ``four \omterm{def:g1:my-body:parts}{limbs}'' claims
        historically: what single thing, long ago, do
        frog, lizard, pigeon, mouse and dolphin all
        descend from?
  \item The class's \omterm{def:g6:classification-kinship:tree}{tree} makes every \omterm{def:g1:plants-around-us:tree}{branch} meet,
        eventually, at the trunk. Say in one sentence
        what that drawing quietly claims about all six
        specimens --- and about the class members drawing
        it.
  \item A skeptic says: ``your \omterm{def:g6:classification-kinship:tree}{tree} is just a filing
        system.'' Give the two-part reply: what the
        nesting of \omterm{def:g6:classification-kinship:attribute}{attributes} is a \emph{fact} about, and
        what the \omterm{def:g6:classification-kinship:tree}{tree} \emph{explains} that mere filing
        would leave unexplained.
\end{enumerate}
\end{problem}
